\documentclass[11pt]{article}
\usepackage[T1]{fontenc}
\usepackage[utf8]{inputenc}
\usepackage{lmodern}
\usepackage[margin=1in]{geometry}
\usepackage{amsmath,amssymb,amsthm,mathtools}
\usepackage{enumitem}
\usepackage[hidelinks]{hyperref}

\newtheorem{theorem}{Theorem}[section]
\newtheorem{definition}[theorem]{Definition}
\newtheorem{lemma}[theorem]{Lemma}
\newtheorem{corollary}[theorem]{Corollary}
\newtheorem{remark}[theorem]{Remark}

\begin{document}

\title{Representatives of Box-Closure Systems on Downsets of Products of Chains}
\author{Luca Blanchi}
\date{June 2026}

\maketitle

\begin{abstract}
We introduce box-closure systems on finite downsets of products of chains. The closure is generated by reflexivity, transitivity, and a coordinate-mixture rule: if two elements lie below the same upper endpoint, then every coordinatewise mixture of them also lies below that endpoint. For a finite downset $D$, we completely characterize all representative sets whose box-closure is the product order on $D$. The classification decomposes by upper endpoint: axis elements have forced incoming cover generators, while every element with at least two active coordinates requires at least one incoming generator from its strict lower ideal. As consequences, inclusion-minimal representatives coincide with cardinality-minimal representatives, all minimum representatives have size $|D|-1$, and after removing the fixed forced axis generators the variable parts of the minimum representatives are precisely the bases of an explicit partition matroid. We derive exact formulae for the number of minimum representatives, Hasse-supported minimum representatives, and all representatives by size, and we solve the weighted representative problem by independent local choices. We also obtain optimal compression ratios for product orders under box-closure and a factorized model of random minimum representatives. A two-dimensional restricted LCA-like system on designated Ferrers cross-pairs is recovered as a motivating special case. The box-closure studied here is not the full LCA $(+)$-closure, nor in general the projection of that closure.
\end{abstract}

\section{Introduction}

Closure systems generated by local inference rules arise naturally in combinatorics, order theory, database theory, and phylogenetics. This paper studies a simple but structured closure system on finite downsets of products of chains. The basic rule is a coordinate-mixture rule: if two lower elements lie below a common upper element, then all coordinatewise mixtures of the two lower elements also lie below the same upper element.

The construction is motivated by restricted cross-consistency rules that appear in LCA-like constraint systems, but the object studied here is not the full LCA $(+)$-closure. This distinction is essential. In the full LCA setting one works with all unordered pairs of leaves, possible non-uniqueness or non-existence of LCAs in directed acyclic graphs, same-side pairs, non-designated cross-pairs, and additional interactions. The projection of a full LCA $(+)$-closure to a designated subsystem need not coincide with the box-closure of that projected subsystem. In this paper, box-closure is treated as an autonomous combinatorial closure operator.

Let
\[
D\subseteq [n_1]\times\cdots\times[n_d]
\]
be a finite downset. For each $x\in D$, introduce a symbol $p_x$. The target relation is the product order
\[
R_D=\{p_x\preceq p_y:x\le y\}.
\]
A representative is a subset $T\subseteq R_D$ whose box-closure is all of $R_D$. The central question is: which subsets $T$ represent $R_D$, and how small can such a representative be?

We solve this problem completely. If $y\in D$ lies on a coordinate axis, then the immediate predecessor relation into $y$ is forced. If $y$ has at least two active coordinates, then any one incoming relation
\[
p_x\preceq p_y,\qquad x<y,
\]
is sufficient. Thus the representative problem decomposes over upper endpoints.

The main consequences are:
\begin{itemize}
\item
\[
\min\{|T|:T^\Box=R_D\}=|D|-1;
\]
\item every inclusion-minimal representative is cardinality-minimal;
\item after removing the fixed forced axis generators, the minimum representatives are bases of a partition matroid;
\item the number of representatives, the number of minimum representatives, and the size enumerator of all representatives have closed product formulae;
\item the weighted representative problem decomposes independently over upper endpoints;
\item for full boxes, the product order has exponentially larger strict relation set than its minimum box-closure representative.
\end{itemize}

The two-dimensional case $D\subseteq [m]\times[n]$ recovers a restricted Ferrers LCA-like system on designated cross-pairs $a_i b_j$. We discuss that connection near the end as motivation, not as the main theorem.

\section{Downsets in products of chains}

Let
\[
[n]=\{1,\dots,n\}
\]
with its usual order. Fix positive integers $n_1,\dots,n_d$, and write
\[
\mathbf n=[n_1]\times\cdots\times[n_d].
\]
We order $\mathbf n$ coordinatewise:
\[
x\le y
\quad\Longleftrightarrow\quad
x_t\le y_t\ \text{for every }t=1,\dots,d.
\]

A subset
\[
D\subseteq \mathbf n
\]
is a \emph{downset} if
\[
y\in D,\quad x\le y
\quad\Longrightarrow\quad
x\in D.
\]
Throughout, $D$ is a nonempty finite downset. Hence $D$ contains the bottom element
\[
\hat{0}=(1,\dots,1).
\]

For $y=(y_1,\dots,y_d)\in D$, define its active support by
\[
\operatorname{supp}(y)=\{t:y_t>1\}.
\]
Then
\[
y=\hat{0}
\quad\Longleftrightarrow\quad
\operatorname{supp}(y)=\varnothing.
\]

A non-minimal element $y\ne\hat{0}$ is called an \emph{axis element} if
\[
|\operatorname{supp}(y)|=1.
\]
If $y$ is an axis element, then it has a unique lower cover, denoted
\[
y^-.
\]
Explicitly, if $t$ is the unique active coordinate, then
\[
y^-=(1,\dots,1,y_t-1,1,\dots,1).
\]

An element $y\in D$ is called \emph{multi-coordinate} if
\[
|\operatorname{supp}(y)|\ge2.
\]

For $y\in D$, write
\[
D_{<y}=\{x\in D:x<y\},
\]
and
\[
D_{\le y}=\{x\in D:x\le y\}.
\]

\section{Box-closure systems}

For each $x\in D$, introduce a formal symbol $p_x$. Let
\[
P_D=\{p_x:x\in D\}.
\]

The target relation is the product order on these symbols:
\[
R_D=\{p_x\preceq p_y:x\le y\}.
\]
Its strict part is
\[
R_D^{<}=\{p_x\preceq p_y:x<y\}.
\]

A generating set will always mean a subset
\[
T\subseteq R_D^{<}.
\]
Reflexive relations are built into the closure and are not counted as generators.

\begin{definition}[Box-closure]
For $T\subseteq R_D^{<}$, the \emph{box-closure} $T^\Box$ is the smallest relation on $P_D$ containing $T$ and satisfying the following rules.

\noindent\textbf{(B0) Reflexivity.}
For every $x\in D$,
\[
p_x\preceq p_x.
\]

\noindent\textbf{(B1) Transitivity.}
If
\[
p_x\preceq p_y
\quad\text{and}\quad
p_y\preceq p_z,
\]
then
\[
p_x\preceq p_z.
\]

\noindent\textbf{(B2) Box rule.}
If
\[
p_x\preceq p_z
\quad\text{and}\quad
p_y\preceq p_z,
\]
then for every coordinatewise mixture $w$ of $x$ and $y$, meaning
\[
w_t\in\{x_t,y_t\}
\qquad(t=1,\dots,d),
\]
with $w\in D$, we add
\[
p_w\preceq p_z.
\]
The box rule preserves the upper endpoint $z$.
\end{definition}

\begin{definition}[Representatives]
A subset $T\subseteq R_D^{<}$ is a \emph{representative} of $R_D$ if
\[
T^\Box=R_D.
\]
The main problem is to characterize all representatives.
\end{definition}

\begin{lemma}[The product order is box-closed]
The relation $R_D$ is closed under the box-closure rules.
\end{lemma}

\begin{proof}
Reflexivity and transitivity are immediate.

Suppose
\[
p_x\preceq p_z
\quad\text{and}\quad
p_y\preceq p_z
\]
belong to $R_D$. Then
\[
x\le z,\qquad y\le z.
\]
Let $w$ be a coordinatewise mixture of $x$ and $y$. For every coordinate $t$,
\[
w_t\in\{x_t,y_t\},
\]
so
\[
w_t\le z_t.
\]
Thus
\[
w\le z.
\]
If $w\in D$, then
\[
p_w\preceq p_z\in R_D.
\]
Hence $R_D$ is box-closed.
\end{proof}

\section{Endpoint preservation}

The first structural observation is that non-reflexive relations cannot be generated without a generator having the same upper endpoint.

\begin{lemma}[Endpoint preservation]
Let $T\subseteq R_D^{<}$. Suppose
\[
p_x\preceq p_y\in T^\Box
\]
with $x<y$. Then $T$ contains at least one generator of the form
\[
p_u\preceq p_y
\]
with $u<y$.
\end{lemma}

\begin{proof}
Let
\[
U(T)=\{y\in D:\exists u<y\text{ such that }p_u\preceq p_y\in T\}.
\]
We prove by induction over derivations in $T^\Box$ that every non-reflexive relation $p_x\preceq p_y$ has $y\in U(T)$.

If the relation is a generator in $T$, the claim is immediate.

Reflexivity produces only relations $p_y\preceq p_y$, so it produces no non-reflexive relation.

For transitivity, suppose
\[
p_x\preceq p_z
\]
is obtained from
\[
p_x\preceq p_y
\quad\text{and}\quad
p_y\preceq p_z.
\]
If $y<z$, then by induction applied to $p_y\preceq p_z$, we have $z\in U(T)$. If $y=z$, then the conclusion is the first premise, and the claim follows from induction applied to that premise.

For the box rule, the conclusion has the same upper endpoint as the two premises. If the conclusion is non-reflexive, then at least one premise is non-reflexive with the same upper endpoint, unless the conclusion coincides with a premise. In either case, induction gives that the common upper endpoint belongs to $U(T)$.

Thus every non-reflexive relation in $T^\Box$ has an upper endpoint that already appears as the upper endpoint of a non-reflexive generator in $T$.
\end{proof}

\begin{corollary}
If $T^\Box=R_D$, then for every $y\ne\hat{0}$, the set $T$ contains at least one generator with upper endpoint $y$.
\end{corollary}

\begin{proof}
For every $y\ne\hat{0}$, choose a lower cover $x\lessdot y$. Then
\[
p_x\preceq p_y\in R_D=T^\Box
\]
is non-reflexive. Apply Lemma 4.1.
\end{proof}

\begin{lemma}[Axis covers are forced]
Let $y$ be an axis element. If
\[
p_{y^-}\preceq p_y\in T^\Box,
\]
then
\[
p_{y^-}\preceq p_y\in T.
\]
\end{lemma}

\begin{proof}
Since $y$ is an axis element, the principal ideal $D_{\le y}$ is a chain:
\[
\hat{0}<\cdots<y^-<y.
\]
In a chain, the box rule cannot create a new lower endpoint: a coordinatewise mixture of two elements of the chain is one of the two elements.

Consider the fiber of relations with upper endpoint $y$. By Lemma 4.1, any non-reflexive relation ending at $y$ ultimately depends on some generator ending at $y$. If that generator is
\[
p_u\preceq p_y
\]
with $u<y^-$, then transitivity can only produce relations
\[
p_z\preceq p_y
\]
with $z\le u$, not the immediate cover relation $p_{y^-}\preceq p_y$. The box rule also cannot create $y^-$ as a new lower endpoint inside this chain.

Therefore the relation
\[
p_{y^-}\preceq p_y
\]
can belong to $T^\Box$ only if it is itself a generator in $T$.
\end{proof}

\section{Complete representative theorem}

We now prove the main classification.

\begin{theorem}[Complete representative classification]
Let $T\subseteq R_D^{<}$. Then
\[
T^\Box=R_D
\]
if and only if the following two conditions hold:
\begin{enumerate}
\item for every axis element $y\ne\hat{0}$,
\[
p_{y^-}\preceq p_y\in T;
\]
\item for every multi-coordinate element $y$, there exists at least one $x\in D_{<y}$ such that
\[
p_x\preceq p_y\in T.
\]
\end{enumerate}
Thus axis elements have forced incoming cover generators, while each multi-coordinate element requires at least one incoming generator from its strict lower ideal.
\end{theorem}

\begin{proof}
We prove necessity and sufficiency.

\noindent\textbf{Necessity.}
Assume
\[
T^\Box=R_D.
\]
By Corollary 4.2, every $y\ne\hat{0}$ occurs as the upper endpoint of at least one generator
\[
p_x\preceq p_y\in T
\]
with $x<y$. This gives condition (2) for all multi-coordinate $y$.

Now let $y$ be an axis element. Since
\[
p_{y^-}\preceq p_y\in R_D=T^\Box,
\]
Lemma 4.3 implies
\[
p_{y^-}\preceq p_y\in T.
\]
Thus condition (1) is necessary.

\noindent\textbf{Sufficiency.}
Assume conditions (1) and (2).

For every $y\ne\hat{0}$, choose a parent $\pi(y)\in D_{<y}$ as follows:
\begin{itemize}
\item if $y$ is an axis element, set
\[
\pi(y)=y^-;
\]
\item if $y$ is multi-coordinate, choose any $x<y$ such that
\[
p_x\preceq p_y\in T.
\]
\end{itemize}
Define
\[
T_\pi=\{p_{\pi(y)}\preceq p_y:y\in D,\ y\ne\hat{0}\}.
\]
Then
\[
T_\pi\subseteq T.
\]
It suffices to prove
\[
T_\pi^\Box=R_D.
\]

We prove by induction on
\[
\rho(y)=y_1+\cdots+y_d
\]
that for every $y\in D$,
\[
p_z\preceq p_y\in T_\pi^\Box
\qquad
\text{for all }z\le y.
\]

If $y=\hat{0}$, this is reflexivity.

First suppose $y$ is an axis element. Then
\[
p_{y^-}\preceq p_y\in T_\pi.
\]
By induction, every $z\le y^-$ satisfies
\[
p_z\preceq p_{y^-}\in T_\pi^\Box.
\]
By transitivity,
\[
p_z\preceq p_y.
\]
Together with reflexivity $p_y\preceq p_y$, this gives every $z\le y$.

Now suppose $y$ is multi-coordinate. Let
\[
x=\pi(y).
\]
The generator
\[
p_x\preceq p_y
\]
belongs to $T_\pi$. Since $x<y$, the induction hypothesis gives
\[
p_{\hat{0}}\preceq p_x.
\]
By transitivity,
\[
p_{\hat{0}}\preceq p_y.
\]

Apply the box rule to
\[
p_{\hat{0}}\preceq p_y
\quad\text{and}\quad
p_y\preceq p_y.
\]
We obtain
\[
p_c\preceq p_y
\]
for every corner $c$ of the box between $\hat{0}$ and $y$, i.e. every $c\in D$ satisfying
\[
c_t\in\{1,y_t\}
\qquad(t=1,\dots,d).
\]

In particular, for every active coordinate $t\in\operatorname{supp}(y)$, let
\[
e_t(y_t)=(1,\dots,1,y_t,1,\dots,1).
\]
Then
\[
p_{e_t(y_t)}\preceq p_y.
\]

Let $z\le y$ be arbitrary. For each coordinate $t$, define
\[
e_t(z_t)=(1,\dots,1,z_t,1,\dots,1).
\]
If $z_t=1$, then $e_t(z_t)=\hat{0}$, and we already know
\[
p_{\hat{0}}\preceq p_y.
\]
If $z_t>1$, then
\[
e_t(z_t)\le e_t(y_t).
\]
By the axis case and transitivity,
\[
p_{e_t(z_t)}\preceq p_y.
\]

Finally,
\[
z=\bigvee_{t\in\operatorname{supp}(z)} e_t(z_t).
\]
Repeated applications of the box rule to the relations
\[
p_{e_t(z_t)}\preceq p_y
\]
produce
\[
p_z\preceq p_y.
\]
Thus the entire principal ideal below $y$ is generated.

The induction is complete. Hence
\[
T_\pi^\Box=R_D.
\]
Since
\[
T_\pi\subseteq T,
\]
we obtain
\[
T^\Box=R_D.
\]
\end{proof}

\section{Minimum and inclusion-minimal representatives}

The complete representative theorem immediately identifies all minimal representatives.

\begin{theorem}[Inclusion-minimal equals cardinality-minimal]
Let $T\subseteq R_D^{<}$ be a representative. Then $T$ is inclusion-minimal if and only if:
\begin{enumerate}
\item for every axis element $y\ne\hat{0}$,
$T$ contains the forced generator
\[
p_{y^-}\preceq p_y;
\]
\item for every multi-coordinate element $y$, $T$ contains exactly one generator
\[
p_x\preceq p_y
\]
with $x<y$.
\end{enumerate}
Consequently, every inclusion-minimal representative has cardinality
\[
|D|-1.
\]
In particular,
\[
\min\{|T|:T^\Box=R_D\}=|D|-1.
\]
\end{theorem}

\begin{proof}
By Theorem 5.1, every representative must contain all forced axis generators and at least one incoming generator for every multi-coordinate element.

If a representative contains two or more incoming generators with the same multi-coordinate upper endpoint $y$, deleting all but one still leaves a representative by Theorem 5.1. Hence an inclusion-minimal representative contains exactly one incoming generator for each multi-coordinate $y$.

Conversely, any set satisfying the two stated conditions is a representative by Theorem 5.1. Removing any generator violates one of the necessary conditions, so the representative is inclusion-minimal.

Such a representative contains exactly one generator for each non-minimal element of $D$. Therefore its cardinality is
\[
|D|-1.
\]
\end{proof}

\section{Partition matroid structure}

Let
\[
F_{\mathrm{ax}}
=
\{p_{y^-}\preceq p_y:
y\ne\hat{0},\ |\operatorname{supp}(y)|=1\}
\]
be the set of forced axis generators.

For every multi-coordinate element $y$, define
\[
E_y=\{p_x\preceq p_y:x<y\}.
\]

Then Theorem 6.1 says that every inclusion-minimal representative is exactly
\[
F_{\mathrm{ax}}\cup \{e_y:y\text{ multi-coordinate},\ e_y\in E_y\}.
\]

\begin{theorem}[Partition matroid theorem]
After removing the fixed forced set $F_{\mathrm{ax}}$, the variable parts of the inclusion-minimal representatives are precisely the bases of the partition matroid
\[
\mathcal M_D=
\bigoplus_{\substack{y\in D\\ |\operatorname{supp}(y)|\ge2}}
U_{1,E_y},
\]
where $U_{1,E_y}$ is the rank-one uniform matroid on $E_y$.
\end{theorem}

\begin{proof}
The variable ground set is
\[
E_{\mathrm{multi}}
=
\bigcup_{\substack{y\in D\\ |\operatorname{supp}(y)|\ge2}} E_y.
\]
The sets $E_y$ are pairwise disjoint because their elements have distinct upper endpoints. By Theorem 6.1, choosing the variable part of an inclusion-minimal representative is exactly choosing one element from each block $E_y$. This is precisely the basis family of the direct sum of the rank-one uniform matroids $U_{1,E_y}$.
\end{proof}

\begin{remark}
Equivalently, one may build a matroid on the full generating set by adding the elements of $F_{\mathrm{ax}}$ as fixed coloops and taking the direct sum with $\mathcal M_D$. We avoid this extra formalism and simply remove the fixed forced set.
\end{remark}

\section{Enumerators and compression}

The preceding results yield exact enumeration formulae and optimal compression ratios.

\subsection{Inclusion-minimal representatives}

The number of inclusion-minimal representatives is
\[
N_{\min}(D)
=
\prod_{\substack{y\in D\\ |\operatorname{supp}(y)|\ge2}}
|D_{<y}|.
\]
Indeed, for each multi-coordinate $y$, one chooses an arbitrary parent $x\in D_{<y}$. Axis generators are forced.

For a full box
\[
D=[n_1]\times\cdots\times[n_d],
\]
we have
\[
|D_{<y}|=\prod_{t=1}^{d} y_t-1.
\]
Therefore
\[
N_{\min}([n_1]\times\cdots\times[n_d])
=
\prod_{\substack{1\le y_t\le n_t\\ |\operatorname{supp}(y)|\ge2}}
\left(\prod_{t=1}^{d} y_t-1\right).
\]

\subsection{Hasse-supported minimum representatives}

A representative is \emph{Hasse-supported} if all its generators are cover relations.

For a multi-coordinate element $y$, the number of lower covers is
\[
|\operatorname{supp}(y)|.
\]
Therefore the number of Hasse-supported minimum representatives is
\[
N_{\mathrm{Hasse}}(D)
=
\prod_{\substack{y\in D\\ |\operatorname{supp}(y)|\ge2}}
|\operatorname{supp}(y)|.
\]

For a full box,
\[
N_{\mathrm{Hasse}}([n_1]\times\cdots\times[n_d])
=
\prod_{\substack{1\le y_t\le n_t\\ |\operatorname{supp}(y)|\ge2}}
|\operatorname{supp}(y)|.
\]

In dimension $2$, this specializes to
\[
N_{\mathrm{Hasse}}([m]\times[n])
=
2^{(m-1)(n-1)}.
\]

\subsection{All representatives}

Theorem 5.1 also counts all representatives.

For an axis element $y$, the forced generator
\[
p_{y^-}\preceq p_y
\]
must be included, while any other incoming generator from
\[
D_{<y}\setminus\{y^-\}
\]
may be included or omitted. Thus the number of choices in that endpoint fiber is
\[
2^{|D_{<y}|-1}.
\]

For a multi-coordinate element $y$, any nonempty subset of $D_{<y}$ may be chosen as the set of incoming generators. Thus the number of choices in that endpoint fiber is
\[
2^{|D_{<y}|}-1.
\]

Therefore the total number of representatives is
\[
N_{\mathrm{all}}(D)
=
\prod_{\substack{y\in D\\ |\operatorname{supp}(y)|=1}}
2^{|D_{<y}|-1}
\cdot
\prod_{\substack{y\in D\\ |\operatorname{supp}(y)|\ge2}}
\left(2^{|D_{<y}|}-1\right).
\]

The ordinary generating function by cardinality is
\[
F_D(t)
=
\prod_{\substack{y\in D\\ |\operatorname{supp}(y)|=1}}
t(1+t)^{|D_{<y}|-1}
\cdot
\prod_{\substack{y\in D\\ |\operatorname{supp}(y)|\ge2}}
\left((1+t)^{|D_{<y}|}-1\right).
\]
The coefficient of $t^k$ is the number of representatives of size $k$.

\subsection{Optimal compression of the product order}

The strict product-order relation has size
\[
|R_D^{<}|
=
\sum_{y\in D}(|D_{\le y}|-1).
\]
The minimum representative has size
\[
|D|-1.
\]
Thus the optimal strict-order compression ratio is
\[
\frac{|R_D^{<}|}{|D|-1}
=
\frac{\sum_{y\in D}(|D_{\le y}|-1)}{|D|-1}.
\]

For a full box
\[
D=[n_1]\times\cdots\times[n_d],
\]
we have
\[
|D|=\prod_{t=1}^{d}n_t
\]
and
\[
|R_D|
=
\prod_{t=1}^{d}\frac{n_t(n_t+1)}2.
\]
Hence
\[
|R_D^{<}|
=
\prod_{t=1}^{d}\frac{n_t(n_t+1)}2
-
\prod_{t=1}^{d}n_t,
\]
while
\[
\min |T|=\prod_{t=1}^{d}n_t-1.
\]

For the cube $D=[n]^d$,
\[
|R_D^{<}|=
\left(\frac{n(n+1)}2\right)^d-n^d,
\qquad
\min |T|=n^d-1.
\]
As $n\to\infty$,
\[
\frac{|R_D^{<}|}{\min |T|}
\sim
\left(\frac n2\right)^d.
\]
Thus box-closure gives an optimal representative that is polynomially smaller, by a factor of order $n^d$, than the full strict product order in fixed dimension $d$.

\section{Weighted optimization and random minimum representatives}

Let a weight
\[
w(x,y)
\]
be assigned to every strict relation
\[
p_x\preceq p_y,\qquad x<y.
\]

\subsection{Minimum-weight inclusion-minimal representatives}

Among inclusion-minimal representatives, the minimum weight is
\[
\sum_{\substack{y\in D\\ |\operatorname{supp}(y)|=1}}
w(y^-,y)
+
\sum_{\substack{y\in D\\ |\operatorname{supp}(y)|\ge2}}
\min_{x<y} w(x,y).
\]

The axis terms are forced, and each multi-coordinate endpoint fiber independently chooses one minimum-weight incoming generator.

\subsection{Minimum-weight representatives with nonnegative weights}

If all weights are nonnegative, then every minimum-weight representative may be chosen inclusion-minimal. Hence the same formula gives the minimum weight among all representatives.

\subsection{Arbitrary real weights}

If negative weights are allowed, extra generators may reduce the total weight. The problem still decomposes by upper endpoint.

For an axis element $y$, the forced generator $p_{y^-}\preceq p_y$ must be included, and any additional negative-weight incoming generator should be included. The optimal axis contribution is
\[
w(y^-,y)
+
\sum_{\substack{x<y\\ x\ne y^-}}
\min\{0,w(x,y)\}.
\]

For a multi-coordinate element $y$, one must choose a nonempty subset of $D_{<y}$. The optimal contribution is
\[
\min_{\varnothing\ne A\subseteq D_{<y}}
\sum_{x\in A} w(x,y).
\]
Equivalently, include all negative-weight incoming generators; if there are no negative-weight incoming generators, include one generator of minimum weight.

Thus the arbitrary-weight representative problem decomposes into independent endpoint-fiber optimizations.

\subsection{Random minimum representatives}

A uniformly random inclusion-minimal representative is sampled as follows:
\begin{enumerate}
\item include every forced axis generator;
\item for each multi-coordinate element $y$, choose one parent $x\in D_{<y}$ uniformly and independently.
\end{enumerate}

Thus the probability that a generator $p_x\preceq p_y$ appears in a uniformly random minimum representative is
\[
\mathbb P(p_x\preceq p_y\in T)
=
\begin{cases}
1, & y\text{ is an axis element and }x=y^-,\\[4pt]
0, & y\text{ is an axis element and }x\ne y^-,\\[4pt]
\dfrac1{|D_{<y}|}, & y\text{ is multi-coordinate and }x<y.
\end{cases}
\]

The entropy of the uniform distribution on minimum representatives is
\[
H_{\min}(D)
=
\sum_{\substack{y\in D\\ |\operatorname{supp}(y)|\ge2}}
\log |D_{<y}|
=
\log N_{\min}(D).
\]

\section{The two-dimensional restricted LCA-like case}

The box-closure system was motivated by a restricted cross-consistency rule for designated LCA-like pairs. We now describe this connection in dimension $2$.

Let
\[
D\subseteq [m]\times[n]
\]
be a Ferrers downset. Write
\[
p_{ij}=a_i b_j.
\]
The designated universe is
\[
P_D=\{a_i b_j:(i,j)\in D\}.
\]

The two-dimensional box rule says that if
\[
p_{i\ell}\preceq p_{km}
\quad\text{and}\quad
p_{rj}\preceq p_{km},
\]
then, whenever the symbols lie in $P_D$, one obtains
\[
p_{ij}\preceq p_{km}
\quad\text{and}\quad
p_{r\ell}\preceq p_{km}.
\]
This is the restricted cross-consistency rule on designated cross-pairs.

This restricted closure is not the full LCA $(+)$-closure on all pairs of leaves. It ignores same-side pairs $a_i a_{i'}$, $b_j b_{j'}$, non-designated cross-pairs, and other interactions present in the full LCA universe. Moreover, the projection of a full LCA $(+)$-closure to the designated universe need not coincide with the box-closure of the designated subsystem.

\subsection{A rooted DAG realizing the designated product order}

For completeness, we recall a realization of the designated product order by a rooted DAG.

For each $(i,j)\in D$, introduce an internal vertex $v_{ij}$. If $(i,j)$ is covered by $(k,\ell)$ in the product order, add a directed edge
\[
v_{k\ell}\to v_{ij}.
\]
If $D$ has several maximal cells, add a root $\rho$ with edges to the maximal vertices. If $D$ has a unique maximal cell, it may serve as the root.

For each represented row $i$, add a leaf $a_i$ below $v_{i1}$. For each represented column $j$, add a leaf $b_j$ below $v_{1j}$.

Then for every designated pair $(i,j)\in D$,
\[
\operatorname{lca}(a_i,b_j)=v_{ij}
\]
is unique in this rooted DAG. Moreover, $v_{ij}$ is a descendant of $v_{k\ell}$ if and only if
\[
i\le k,\qquad j\le \ell.
\]
Thus the LCA order induced on the designated cross-pairs
\[
P_D=\{a_i b_j:(i,j)\in D\}
\]
is exactly the product order $R_D$.

This statement is only about designated cross-pairs. The full LCA relation of the DAG may contain additional comparisons involving other pairs of leaves. If one wants the DAG to satisfy a particular convention for phylogenetic networks, degree-$(1,1)$ vertices may be suppressed or the construction may be modified in the standard way. The representative theorems above do not depend on these network conventions.

\subsection{Two-dimensional representatives}

The two-dimensional results are immediate specializations of the general theory.

The axis elements are the first row and first column. Their immediate predecessor generators are forced. Every interior cell $(i,j)$, $i,j>1$, requires exactly one incoming generator from its strict lower ideal in an inclusion-minimal representative.

Thus
\[
\min\{|T|:T^\Box=R_D\}=|D|-1.
\]

For a rectangle
\[
D=[m]\times[n],
\]
the number of inclusion-minimal representatives is
\[
\prod_{i=2}^{m}\prod_{j=2}^{n}(ij-1),
\]
and the number of Hasse-supported minimum representatives is
\[
2^{(m-1)(n-1)}.
\]

\section{Relationship with LCA constraints and triple representatives}

The classical LCA-constraint problem concerns rooted trees and constraints comparing lowest common ancestors. Recent work extends LCA constraints to directed acyclic graphs and phylogenetic networks, where LCA existence and uniqueness become nontrivial and where closure operations involve the full universe of leaf pairs.

The present paper studies a different restricted closure system. Its relation to LCA theory is motivational and occurs explicitly in the two-dimensional designated cross-pair case. The main results should not be read as theorems about the full LCA $(+)$-closure.

Representative triple sets in rooted phylogenetic trees have a known matroidal structure. The box-closure systems studied here also yield a matroidal structure, but of a different and simpler kind: after removing the forced axis generators, the variable parts of the inclusion-minimal representatives form the bases of a partition matroid. This follows from endpoint-fiber decomposition, not from triple-closure theory.

\section{Conclusion}

We introduced box-closure systems on finite downsets of products of chains and solved their representative problem completely.

The main results are:
\begin{enumerate}
\item a complete characterization of all representatives;
\item the equivalence of inclusion-minimality and cardinality-minimality;
\item the minimum size formula $|D|-1$;
\item a partition-matroid description of the variable parts of minimum representatives after removing forced axis generators;
\item exact enumerators for minimum, Hasse-supported, and all representatives;
\item optimal compression ratios for product orders under box-closure;
\item a decomposed solution of the weighted representative problem;
\item a factorized model of random minimum representatives;
\item recovery of the two-dimensional restricted Ferrers LCA-like system as a special case.
\end{enumerate}

The contribution is a complete combinatorial analysis of a coordinate-mixture closure system. Its connection to LCA constraints is motivational and restricted, not a claim about the full LCA $(+)$-closure.

\begin{thebibliography}{99}
\small
\sloppy

\bibitem[ASSU81]{ASSU81}
A. V. Aho, Y. Sagiv, T. G. Szymanski, and J. D. Ullman,
\emph{Inferring a tree from lowest common ancestors with an application to the optimization of relational expressions},
SIAM Journal on Computing 10 (1981), 405--421.

\bibitem[EH26]{EH26}
P. A. Ebert and M. Hellmuth,
\emph{Inferring phylogenetic networks from required and forbidden LCA-constraints},
arXiv:2605.03827, 2026.

\bibitem[HLM26]{HLM26}
M. Hellmuth, A. Lindeberg, and V. Moulton,
\emph{Encoding phylogenetic networks with least common ancestor constraints},
arXiv:2606.16963, 2026.

\bibitem[HS18]{HS18}
M. Hellmuth and C. R. Seemann,
\emph{The matroid structure of representative triple sets and triple-closure computation},
European Journal of Combinatorics 70 (2018), 384--407.

\bibitem[LAMSH25]{LAMSH25}
A. Lindeberg, A. Alfonsson, V. Moulton, G. E. Scholz, and M. Hellmuth,
\emph{Inferring DAGs and phylogenetic networks from least common ancestors},
arXiv:2511.07965, 2025.

\end{thebibliography}

\section*{Declaration of generative AI and AI-assisted technologies in the writing process}
During the preparation of this work, ChatGPT, by OpenAI, was used to assist with mathematical drafting, formalization, review, and editing. This work is shared as a preliminary AI-assisted mathematical note. The mathematical content may have been only partially reviewed and may contain errors; it should not be treated as peer-reviewed or as a fully verified manuscript.

\end{document}
